% results.tex -- Results
\section{Results}

To explore how investor decisions could destabilize a market, we changed how
some traders responded to their neighbors and found that a majority entered
the Sell state only when enough traders adopted the new decision rule
(Fig.~\ref{fig:rule-quench}; Supplementary Table~\ref{tab:paired-switch}). We compared each market with an otherwise
identical market where traders continued to use their original decision
rules. The new rule, Rule~19286, changed the response to three neighbors buying
and one holding from Buy to Sell, while keeping every other response
unchanged (Fig.~\ref{fig:rule-quench}A,B). We assigned it to four different
fractions of traders inside the grid. We
used a $64\times64$ grid with an unlimited dealer and technical traders along
its edge who decided whether to buy, hold, or sell from the price
(Supplementary Table~\ref{tab:campaigns}). Each pair of markets had the same
1,000-step history and the same price immediately after we assigned the new
rule. They first differed when the selected traders next decided how to act
with three neighbors buying and one holding. All 50 markets that kept the
original rule remained below a Sell majority for the next 3,000 steps. The
same was true when we changed the rule for a fraction $f=0.10$ or 0.15 of
interior traders, whereas changing it for $f=0.20$ or 0.25 produced a
majority in every run. The median time to reach a majority fell from 1,044.5
to 311.5 steps between the two larger fractions
(Supplementary Table~\ref{tab:paired-switch}).
In the representative $f=0.20$ pair, groups of sellers expanded across the
grid after the rule change while decisions in the comparison market remained
mixed (Fig.~\ref{fig:rule-quench}C). We also tested whether a temporary
decision to sell had the same effect as changing the rule for later
decisions. Setting the same selected traders to Sell once, while keeping
their original rule, produced no majority at $f=0.15$, 0.20, or 0.25,
whether the technical traders were active or kept in Hold (Supplementary
Table~\ref{tab:revision-pulse}). Of all 18 ways to change a single response
in the original rule, Rule~19280, six produced a majority at $f=0.20$ and
the other twelve did not (Supplementary Table~\ref{tab:revision-rules}).
Within these tests, the shift to a market dominated by sellers depended on
how many traders changed their decision rule, which decision changed, and
whether the new rule remained in use.

Traders inside the grid decided from their neighbors' actions, while the
technical traders along its edge decided from the price and could therefore
pass the effects of a price change to their neighbors. To determine whether a Sell majority could
form without that influence, we changed the decision rule for selected
interior traders while keeping the technical traders in Hold, and found
that a majority still formed (Fig.~\ref{fig:reservoir}). With the technical
traders unable to buy or sell, the fraction of runs reaching a majority rose
from 0.34 at $f=0.175$ to 0.70 at 0.185 and 1.00 at 0.20
(Fig.~\ref{fig:reservoir}A). Changing dealer capacity left the entire
sequence of interior decisions unchanged for each starting market and
fraction $f$, even though completed trades, prices, and dealer holdings
differed. Prices could still change after trades, but no trader could use those
prices to change a decision. Responses to neighboring traders were therefore
sufficient to sustain widespread selling. Allowing the technical traders to
trade on the price again changed how often a majority formed only among
fractions where some runs reached a majority and others did not
(Fig.~\ref{fig:reservoir}A,B). With these traders active and an unlimited
dealer, 14 of 50 runs reached a majority at $f=0.175$, 43 did so at 0.185,
and all 50 did so at every tested fraction from 0.195 upward. Across all
dealer capacities with these traders active, no run reached a majority at
$f\leq0.15$ and every run did so at $f\geq0.20$. Responses to price thus
changed how often a majority formed but were not required for it to develop.
We next asked how this instability depended on the number of interacting
traders and the time available, and when it was accompanied by a large price
decline.

The fraction of traders needing the new rule could depend on how many
traders interacted and how much time their decisions had to influence
one another. We changed the rule at five grid sizes while keeping the technical
traders in Hold, and found that larger markets needed a smaller fraction
of traders to receive the new rule for a majority to form
(Fig.~\ref{fig:finite-size}). We tested additional fractions near the
rise in majority probability to measure its position and width. After
3,000 steps, the fraction at which half the runs reached a majority,
$f_{50}$, fell from 0.216 at $L=32$ to 0.166 at
$L=128$. The additional fraction needed to raise majority probability
from 0.25 to 0.75 decreased from 0.023 [0.018, 0.027] to
0.008 [0.004, 0.010], with brackets giving 95\% confidence
intervals obtained by resampling the starting markets
(Fig.~\ref{fig:finite-size}A,C; Supplementary Table~\ref{tab:finite-size-scaling}).
We also measured variation between runs by multiplying the variance of
the final mean Sell fraction by the number of interior traders. The maximum value of this measure, susceptibility, rose from
21.2 to 241.2 over the
same sizes (Fig.~\ref{fig:finite-size}B,D). Fitting smooth curves that
required probability to increase with rule adoption placed each
midpoint within 0.004 of the estimate from interpolation between tested
fractions, although the fitted widths were narrower when the technical traders
stayed in Hold (Supplementary Table~\ref{tab:revision-fits}). Allowing
more time also changed the result: extending the $L=64$ runs from 3,000
to 9,000 steps lowered the fraction needed for half the runs to reach a
majority from about 0.178 to 0.167. The larger of the two endpoint grids still needed a lower fraction
when we adjusted observation time with grid size or changed the Sell
fraction required to count a transition, but the estimated midpoints
kept moving in longer runs
(Supplementary Table~\ref{tab:revision-horizons}). The sharper rise in
majority probability on larger grids therefore described the markets
over the time observed; it did not establish a threshold that would
remain fixed as observation time increased.


To determine how changes in investor decisions developed into market
instability, we followed five selected markets beyond the original
observation period and found that groups of sellers continued to grow
(Fig.~\ref{fig:dose-panels}). We retained five examples selected from their
original 3,000-step results before inspecting any images, then extended each
to 10,000 steps after the rule change. The displayed $f=0.15$ market retained
a mixture of decisions throughout this longer period. Its price remained
roughly steady after an initial decline, with maximum drawdown increasing
only from 5.3\% to 5.4\% (Supplementary
Table~\ref{tab:market-panel-horizons}). The two $f=0.175$ markets both
reached a Sell majority, but at different times. One reached a majority
after 2,186 steps; the other remained below one-half at the original endpoint
and reached a majority after 4,117 steps. The $f=0.185$ and $f=0.20$ markets
reached a majority after 1,779 and 1,030 steps, respectively
(Fig.~\ref{fig:dose-panels}A,C). In each of these four markets, groups of
sellers enlarged and joined as the price declined. In the later $f=0.175$
transition, maximum drawdown grew from 9.9\% at 3,000 steps to 25.6\% at
10,000 steps. The other three markets ended with drawdowns of 26.5\%,
26.3\%, and 26.4\% (Supplementary Table~\ref{tab:market-panel-horizons}).
A market that had not reached a Sell majority by the original endpoint
could therefore develop a majority and a large price decline when observed
for longer. We also extended all 50 original $f=0.15$ starting markets and
their matched controls. Two markets with the new rule reached a majority
by 10,000 steps, compared with none at 3,000 steps, while all controls
remained below a majority. Median maximum drawdown in the changed-rule
markets rose from 4.9\% to 5.5\%, but the largest increased from 9.3\% to
24.5\% (Supplementary Table~\ref{tab:market-panel-horizons}). The displayed
$f=0.15$ example thus illustrated persistent mixed decisions over this
period, while other starting markets at the same fraction developed delayed
instability. The growing groups connected individual decisions to a market
dominated by sellers, and the longer observations showed that delayed
market instability could accompany a large price decline.

A decision to sell could influence neighboring traders without finding a
buyer, so a market dominated by sellers did not necessarily undergo a price
crash. We varied the dealer's ability to buy and sell across eight
capacities with the technical traders active, and found
that greater capacity increased the price decline but did not consistently
increase the probability of a Sell majority (Fig.~\ref{fig:reservoir}). At
$f=0.175$, the probability was 0.66 with capacity $h=0.1$ and 0.28 with
unlimited capacity, a difference of 38 percentage points in favor of the
smaller dealer (Fig.~\ref{fig:reservoir}A,B). The technical traders normally
adjusted the width of their band, the departure from the reference price
that triggered a decision to buy or sell. Keeping that required departure fixed at its value
before the rule change changed the probability difference from 0.38 to
$-0.02$. This linked the smaller dealer's advantage to traders adjusting
their responses to price, although other increases and decreases in
probability across capacities remained (Supplementary
Table~\ref{tab:revision-band}). At $f=0.20$, every run reached a majority at
every capacity, but a larger dealer allowed more orders backed by trader
cash or shares to complete. Median maximum drawdown rose from zero at $h=0$
to 0.40\% at $h=0.1$, 3.92\% at $h=1$, and 25.3\% with unlimited capacity,
while the median fraction of these orders rejected by the dealer fell from
0.449 to zero (Fig.~\ref{fig:reservoir}C,D). Without dealer capacity,
traders could only buy from or sell to one another, so equal numbers of
completed buys and sells left the price unchanged. Capacity $h=10$ rejected
no orders and matched unlimited capacity for the resources and duration
tested. Larger dealers still produced larger drawdowns when we varied how
strongly completed trades changed the price over a tenfold range, although
drawdown magnitudes changed by a factor of seven to ten (Supplementary
Table~\ref{tab:revision-impact}). Runs from the same starting market also
separated selling decisions from price declines. At $f=0.20$, the market
initialized with seed 25 reached a majority at $h=0.1$, 1, and $\infty$
within 88 steps of one another. The mean Sell fraction over the final 500
steps lay between 0.81 and 0.87 in these three runs and in the comparison
with the technical traders kept in Hold, yet maximum drawdown rose from 0.40\% at
$h=0.1$ to 3.92\% at $h=1$ and 25.8\% with unlimited capacity
(Fig.~\ref{fig:dealer-panels}). Similar instability in decisions thus
accompanied both a small decline and a simulated price crash. Dealer
capacity limited the price decline after most interior traders had decided
to sell, while the resulting price influenced later decisions.

Traders needed shares to carry out a decision to sell and cash to carry out
a decision to buy, even when their decision rules called for those actions.
To distinguish the effects of trader resources from the dealer's ability to
complete their orders, we multiplied trader cash and shares independently by
0.1, 1, or 10 when we changed the rule, and found that holdings changed
majority probability only when the technical traders were active
(Fig.~\ref{fig:balance-sheet}). We tested these changes at three dealer
capacities, both with the technical traders active and with them kept in
Hold. At $f=0.18$, all nine combinations of cash and shares
produced identical interior decisions when the technical traders stayed in Hold, with
a majority probability of 0.64 across dealer capacities. When the technical traders
were active, probabilities ranged from 0.50 to 0.94, and the other
tested fractions showed the same dependence on whether these traders were
active (Fig.~\ref{fig:balance-sheet}A; Supplementary
Table~\ref{tab:balance-fractions}). Holdings determined whether traders
could carry out their decisions, but changed later decisions only through
their effect on price. With the original cash holdings and an unlimited
dealer, increasing the share multiplier from 0.1 to 1 to 10 raised median
drawdown from 2.87\% to 20.1\% to 83.5\%, because every sell order backed by
a share completed and the dealer bought all shares that other traders did
not buy (Fig.~\ref{fig:balance-sheet}B). Limited dealer capacity capped that
increase: the original and tenfold share holdings both produced 3.92\% at
$h=1$ and 0.399\% at $h=0.1$, because the dealer rejected sell orders beyond
its capacity (Fig.~\ref{fig:balance-sheet}C). Majority probability did not
follow drawdown (Fig.~\ref{fig:balance-sheet}D). At $h=1$ with the original
share holdings, tenfold cash reduced median drawdown from 3.92\% to 0.089\%
while raising majority probability from 0.78 to 0.94. The extra cash allowed
more traders to buy, so they bought shares that would otherwise have needed
the dealer as a buyer, while the changed price also influenced the technical
traders' decisions. Trader and dealer resources therefore determined whether a market
dominated by sellers experienced a limited decline or a large price crash.
The decision to sell and the resources needed to complete the sale
contributed separately to the market outcome.
